\documentclass[11pt,a4paper]{article}
\usepackage[T1]{fontenc}
\usepackage[utf8]{inputenc}
\usepackage{lmodern}
\usepackage[a4paper,margin=26mm,headheight=15pt]{geometry}
\usepackage{microtype,amsmath,amssymb,amsthm,mathtools}
\usepackage{booktabs,longtable,array,tabularx}
\usepackage{xcolor,listings,enumitem}
\usepackage{fancyhdr}
\usepackage{xurl}
\usepackage[unicode,colorlinks=true,linkcolor=blue!45!black,citecolor=blue!45!black,urlcolor=blue!45!black]{hyperref}
\usepackage[nameinlink,noabbrev]{cleveref}
\usepackage{bookmark}
\hypersetup{pdftitle={Forge over Infinite Ordered State Spaces},pdfauthor={Research prepared for Vladimir Reshetnikov},pdfsubject={Exact unsafe-region certificates for resources and lossy queues},pdfkeywords={Lean, Forge, coverability, well-quasi-orders, certificates, lossy channels}}
\setlength{\parskip}{4pt}
\setlength{\parindent}{1.2em}
\setlength{\emergencystretch}{2em}
\setcounter{tocdepth}{2}
\pagestyle{fancy}
\fancyhf{}
\fancyhead[L]{\small\textsc{Forge / Ordered-state certificates}}
\fancyhead[R]{\small Research and implementation note}
\fancyfoot[C]{\thepage}
\renewcommand{\headrulewidth}{0.3pt}
\lstset{basicstyle=\small\ttfamily,columns=fullflexible,keepspaces=true,
  breaklines=true,showstringspaces=false,frame=single,framerule=0.3pt,
  rulecolor=\color{black!25},xleftmargin=0.5em,xrightmargin=0.5em,
  aboveskip=8pt,belowskip=8pt,tabsize=2}
\newtheorem{theorem}{Theorem}[section]
\newtheorem{lemma}[theorem]{Lemma}
\newtheorem{proposition}[theorem]{Proposition}
\newtheorem{corollary}[theorem]{Corollary}
\theoremstyle{definition}
\newtheorem{definition}[theorem]{Definition}
\newtheorem{example}[theorem]{Example}
\newtheorem{remark}[theorem]{Remark}
\newcommand{\N}{\mathbb N}
\newcommand{\Up}[1]{\mathord{\uparrow}#1}
\newcommand{\Pre}{\operatorname{Pre}}
\newcommand{\Bad}{\mathsf{Bad}}
\newcommand{\Unsafe}{\mathsf{Unsafe}}
\newcommand{\Safe}{\mathsf{Safe}}
\newcommand{\Min}{\operatorname{Min}}
\newcommand{\Reach}{\operatorname{Reach}}
\newcommand{\eps}{\varepsilon}
\newcommand{\vct}[1]{\boldsymbol{#1}}
\newcommand{\code}[1]{\texttt{#1}}
\newcommand{\status}[1]{\noindent\textbf{Implementation status.} #1\par}
\newcommand{\proposed}{\textit{Proposed interface; not an implemented Lean API.}}
\title{\vspace{-1.5em}\textbf{Forge over Infinite Ordered State Spaces}\\[0.7em]
\large Exact unsafe-region certificates for resources and lossy queues\\[0.3em]
\normalsize A non-duplicating extension to the reviewed Forge design}
\author{Research and prototype prepared for Vladimir Reshetnikov}
\date{15 September 2026\\\small Repository snapshots and limitations are specified in Section~1}
\begin{document}
\maketitle
\thispagestyle{empty}
\begin{abstract}
This article proposes two specialized workers for Forge: one for guarded
nonnegative-affine resource transitions, including ordinary token consumption,
reset, transfer and duplication; the other for finite-control lossy FIFO
channels. Both operate over genuinely infinite state spaces. Their result is
not merely a verdict about one initial state, but an exact finite description
of \emph{every} initial state from which a finite bad execution is possible.

The mathematical foundation is classical backward coverability in
well-quasi-ordered systems. The contribution is a concrete, independently
checkable implementation design for Forge: an immutable derivation DAG supplies
one inclusion, and complete predecessor-coverage obligations supply the other.
For affine resources, a proved clipping bound reduces infinite predecessor
sets to integer boxes. For lossy queues, exact word-predecessor rules preserve
FIFO order while accounting explicitly for message loss. A second producer
checks closure directly against the current global antichain, avoiding the
construction of large, irrelevant intermediate Pareto bases.

The accompanying standard-library Python implementation was executed. Its
first producer emitted 190 accepted region certificates. Local differential
tests made 63,504 resource-point comparisons and 2,325 word-pair comparisons.
Complete finite execution oracles agreed on 7,000 resource queries, 2,040
acyclic-channel queries and 2,480 length-nonincreasing channel queries. A
separation-based producer returned the identical bases on the same 190 models.
On three deliberately constructed dense-preimage cases it completed where the
eager producer reached its stated grid limit. One of these avoids explicitly
listing 26,075,972,546 local minima; its full compact-JSON certificate is 1,063
bytes.

These are Python experiments, not Lean-kernel proofs. No Lean executable was
available for this work, no Lean tactic was implemented, and no comparison
against \code{grind} was run. The article specifies the missing Lean theorems,
source-model bridges, integration points and acceptance tests rather than
presenting that work as complete.
\end{abstract}
\clearpage
\tableofcontents
\clearpage

\section{The extension boundary: what this adds, and what it does not}
\label{sec:delta}

\subsection{Reviewed snapshots}
The primary reference is Forge at commit
\begin{center}\small\path{58ea206bd7ad501b930add568151217bcc7f82a2}.\end{center}
The review used its repository tree, root README, implementation-status
inventory, Leant/Djex integration review, and merged finite-closure section.
The auxiliary projects were inspected at Leant
\path{6bf05ad78c467989e68290f2d08bbed40802d485} and Djex
\path{e8778f4ebd63e1f9b9b410fa4de8d14a8a04c9e5}; the current indexing diagnostics
were also available in their commit records. The precise URLs and audit scope
are retained in \path{docs/novelty-audit.md} and
\path{docs/source-manifest.json}.\cite{forge,forgestatus,forgeclosure,leant,djex}

This was \emph{not} a successful full local clone followed by an exhaustive
review of every archived submission. Network access prevented the clone;
connector reads provided the cited source material. A GitHub search returned
an incomplete-results flag, so its zero matches were not used as proof of
absence. Accordingly, ``new'' below means a concrete capability not present in
the reviewed merged capability inventory and closure discussion, not a claim
of priority over all literature or every archived file.

The distinction matters because the current Forge repository already contains
three rounds of substantial work. It describes structural proof planning,
observable-space and polynomial-ideal closure, finite datatype covers,
noncommutative certificates, analytic estimates, quantitative probabilistic
reasoning, and pushdown summaries. It also already establishes the intended
separation between a producer and an acceptance boundary. Those ideas are
prerequisites here, not new contributions.\cite{forgestatus,forgeclosure}

\subsection{A narrowly defined capability delta}
\begin{table}[htbp]
\centering\small
\begin{tabularx}{\textwidth}{@{}p{0.28\textwidth}XX@{}}
\toprule
Reviewed Forge capability & Reused, not reproposed & Added here\\
\midrule
Linear and polynomial invariant closure & Target-directed finite certificates & Exact backward unsafe sets for ordered integer and word states\\
Finite-state and datatype covers & Finite records describing unbounded behavior & Dickson/Higman termination with unbounded counters and FIFO lengths\\
Pushdown summaries & Finite-control reachability vocabulary & FIFO channels with losses, not LIFO stacks\\
Polynomial cone and box certificates & Arithmetic leaves and coverage discipline & Nonnegative-affine integer preimage clipping and direct global coverage\\
Counterexample records and witness synthesis & Explicitly replayable evidence & A finite family of action words covering \emph{all} unsafe initial states\\
Leant/Djex evidence boundaries & Exact subjects and approximation polarity & Model-specific simulation obligations for counters, queues and population quotients\\
\bottomrule
\end{tabularx}
\caption{The contribution is two executable semantic workers and their particular certificates, not another generic controller.}
\end{table}

A semaphore theorem can be proved by a familiar hand-written invariant. That
is not the strongest reason to add this worker. The stronger reason is that a
single checked basis can classify \emph{all} natural initial populations and
identify exactly which extra resources permit a bad execution. Similarly,
queue reasoning retains an unbounded amount of ordered data without requiring
a finite queue-length cutoff. Neither feature requires the source user to
supply the final invariant.

The new closure-separation producer also addresses an integration issue that
a bare instruction to ``use a Petri-net solver'' would miss: a local predecessor
basis can be enormous even when the existing global unsafe region covers it
with one cone. The certificate language should be able to express the cheap
global proof without materializing that intermediate basis.

\subsection{Relationship to established research}
Backward coverability based on well-quasi-orderings is classical, notably in
the work of Abdulla, \v{C}er\=ans, Jonsson and Tsay, and Finkel and
Schnoebelen.\cite{acjt,fs} The lossy-word predecessor operations below are
single-letter specializations of established channel-machine
operations.\cite{flat} This article does not rename those results as new
mathematics. Its original deliverables are the selected Forge extension,
explicit certificate formats, two working producers, a separately written
checker, concrete witnesses, tests, and a port plan that preserves the exact
meaning of the original Lean obligation.

\section{One exact question over two infinite domains}
\label{sec:semantics}

\subsection{Unsafe means existential finite reachability}
Let $S$ be a state space, $T$ a finite set of transitions, and $\to_t$ the
relation for transition $t$. Write $\to=\bigcup_{t\in T}\to_t$ and let
$\to^*$ be its reflexive transitive closure. For a designated bad set,
\begin{equation}
 \Unsafe = \{s\in S:\exists z\in\Bad,\ s\to^*z\},
 \qquad \Safe=S\setminus\Unsafe.
 \label{eq:unsafe}
\end{equation}
Thus $s\in\Safe$ means that \emph{every finite execution} from $s$ avoids
$\Bad$. Membership in $\Unsafe$ means that \emph{some} bad execution exists.
It does not mean all executions fail, that failure is inevitable under a fair
scheduler, or that a controller wins a game. No probabilities are involved.

The user may ask about one initial state, a symbolic family $I(\theta)$, or
an initial predicate $\mathcal I$. The worker computes the same reusable
$\Unsafe$ once. The remaining proof is respectively a finite membership test,
a parameter formula, or
\begin{equation}
 \forall s,\quad \mathcal I(s)\Longrightarrow s\notin\Unsafe.
 \label{eq:initcheck}
\end{equation}

\subsection{Order and upward bases}
A quasi-order $\preceq$ is reflexive and transitive. It is a
\emph{well-quasi-order} (WQO) if every infinite sequence $s_0,s_1,\ldots$
contains $i<j$ with $s_i\preceq s_j$. For $B\subseteq S$, define
\[
 \Up B=\{s:\exists b\in B,\ b\preceq s\}.
\]
We represent the bad set by a finite input list $U$ with
$\Bad=\Up U$. For the two partial orders below, a finite irredundant basis is
an antichain: no two distinct members are comparable.

States have a finite control component $q\in Q$. Order comparisons require
\emph{equal} controls; a state at an error location must not dominate a state
at a different program point merely because its data are smaller. The data
order is componentwise natural order for resources and componentwise
subsequence order for queue contents.

\begin{definition}[Strong upward compatibility]
A transition relation is strongly upward compatible if
\[
 s\preceq s'\ \wedge\ s\to_t z
 \quad\Longrightarrow\quad
 \exists z',\ s'\to_t z'\ \wedge\ z\preceq z'.
 \label{eq:monotone}
\]
The same transition label is used on both sides.
\end{definition}

Both supported languages satisfy this condition. Consequently the predecessor
of an upward set is upward. The entire unsafe set is upward too: induct on
an execution from the smaller state and lift each transition.

\subsection{Why upward targets are an explicit restriction}
``At least two processes are critical'' is upward in resource counts.
``The control is error'' is upward in queue contents. ``The final queue
contains $ab$ as a subsequence'' is also upward. In contrast, an exact counter
value, an exact queue word, or a queue being empty at an otherwise unrestricted
control is not an upward property.

Replacing an exact target $u$ by $\Up{\{u\}}$ changes the question from
reachability to coverability. It can be a deliberately certified
\emph{over-approximation} for a safety proof, but it cannot silently count as
an exact translation. In particular, representing a bad-control target by the
empty word is valid because \emph{all} queues at that control are bad, not
because arbitrary empty-queue reachability has been solved.

\section{A finite certificate for an exact infinite region}
\label{sec:certificate}

\subsection{Two inclusions, with different evidence}
The worker proposes a final basis $B$. To establish
\begin{equation}
 \Unsafe=\Up B,
 \label{eq:exact}
\end{equation}
it must establish both inclusions. These are logically independent.

For $\Up B\subseteq\Unsafe$, it gives a finite acyclic derivation graph.
A node is a concrete state, together with one of two reasons:
\begin{enumerate}[nosep,leftmargin=2em]
\item It is exactly one of the original bad basis states $u\in U$.
\item It names a transition $t$ and an \emph{earlier} node $v$, and its state
      $s$ can take a $t$-step to some $z\succeq v$.
\end{enumerate}
The final basis refers to nodes in this immutable graph. It is legal for a
node to disappear from the active antichain while remaining in the derivation
graph: later proofs may still depend on it.

For $\Unsafe\subseteq\Up B$, it gives target coverage and backward closure:
\begin{align}
 \Up U&\subseteq\Up B,\label{eq:targetcover}\\
 \Pre_t(\Up{\{b\}})&\subseteq\Up B
 &&\text{for every }b\in B\text{ and every }t\in T.
 \label{eq:closurecover}
\end{align}
Control-incompatible pairs have empty predecessor sets and are discharged by
checking their controls. The checker, not the producer, determines the full
set of remaining transition/basis pairs.

\begin{theorem}[Exact-region certificate soundness]
\label{thm:exact}
Assume strong upward compatibility, a finite bad basis, valid derivation
nodes for every member of $B$, and \eqref{eq:targetcover}--\eqref{eq:closurecover}.
Then \eqref{eq:exact} holds. No WQO assumption is required for this theorem.
\end{theorem}
\begin{proof}
First prove by induction on node index that every state above that node can
reach $\Bad$. At a bad leaf this follows from upward closure. At a predecessor
node $s$, the recorded step reaches above the earlier child. From any
$s'\succeq s$, compatibility lifts that step to another state above the same
child; apply the induction hypothesis. Every state in $\Up B$ is above one
of these nodes, proving the first inclusion.

For the second inclusion, induct backwards on the length of a concrete bad
execution. Its final state lies in $\Up U\subseteq\Up B$. If $x\to_t y$ and
$y\in\Up B$, choose $b\in B$ with $b\preceq y$. Then
$x\in\Pre_t(\Up{\{b\}})$, and the corresponding closure obligation puts
$x$ in $\Up B$. This reaches the initial state of every finite bad execution.
\end{proof}

\begin{remark}[Why the two halves must not be conflated]
Backward closure alone accepts the whole state space and says nothing about
whether its states really can fail. Derivations alone describe only discovered
failures and cannot justify safety of an unrepresented state. Their
combination proves exactness. A budget cutoff may still leave useful verified
witnesses, but it does not manufacture the missing closure half.
\end{remark}

\subsection{The eager backward algorithm}
For each principal cone, suppose an exact finite predecessor basis is
computable:
\[
 \Pre_t(\Up{\{b\}})=\Up{P_t(b)}.
\]
The baseline producer uses the following worklist. It is included as an
executable reference and an ablation baseline, not as a new coverability
algorithm.\cite{acjt,fs}
\begin{lstlisting}
active := empty antichain
nodes  := immutable vector of derivations
queue  := empty
admit every original bad target

while queue is not empty:
    n := pop_front(queue)
    if n is no longer active: continue
    for each original transition ending at control(n):
        for p in exact_predecessor_basis(transition, state(n)):
            if some active b <= p: continue
            append node(p, transition, child=n)
            remove active states >= p, retaining their graph nodes
            activate the new node and enqueue it

rebuild all final transition/basis coverage obligations
return nodes, final active basis, target coverage, closure receipts
\end{lstlisting}

\begin{theorem}[Termination and completeness of the idealized producer]
\label{thm:termination}
If $S$ is WQO, every local predecessor computation terminates with an exact
finite basis, and no operational cutoff is imposed, this worklist terminates
and produces a certificate for \eqref{eq:exact}.
\end{theorem}
\begin{proof}
The upward set represented by the active antichain never shrinks. Removing
an active $v$ because a newly admitted $p\preceq v$ preserves all states it
covered. If infinitely many states $s_0,s_1,\ldots$ were admitted, WQO would
give $i<j$ with $s_i\preceq s_j$. At admission time of $s_j$, the current
upward region still contains $s_i$, and hence contains $s_j$. That contradicts
the admission test. Thus only finitely many nodes are admitted. Each is
popped once, each transition list is finite, and each local computation
terminates, so the worklist finishes.

Every admitted node has the recorded valid derivation. Target coverage is
preserved from initialization. At termination, every surviving node was
expanded. Its predecessors were either admitted or already covered; later
changes only enlarge the represented upward set. Therefore all final
closure obligations hold, and \cref{thm:exact} applies.
\end{proof}

The actual program has limits on nodes, expansions and arithmetic partition
work. Its completed finite receipt is checked without trusting this termination
argument. Exhausting a limit returns \code{unknown}, not \code{safe} and not
``the theorem is false.''

\subsection{Minimality and finite witness plans}
For our partial orders, exactness and antichain irredundancy imply that the
final basis is precisely the set of minimal unsafe states. To see this, let
$b\in B$ and suppose an unsafe $x\prec b$ existed. Exactness supplies
$b'\in B$ with $b'\preceq x\prec b$, contradicting the antichain property.
Conversely every minimal unsafe state lies above some basis member and must
equal it.

\begin{corollary}[A finite family of universal failure plans]
\label{cor:plans}
Each final basis node determines a finite action word that works from every
state above that node, with suitable concrete states and, in the FIFO case,
suitable losses. If the certificate contains $N$ nodes, these words have
length at most $N-1$.
\end{corollary}
\begin{proof}
Follow the strictly decreasing child indices to a bad leaf. Compatibility
lifts the same transition labels from any larger starting state. A strictly
decreasing chain of indices uses at most $N-1$ edges.
\end{proof}
This is a bound on the supplied witness family, not a claim that the extracted
word is globally shortest. It is also not a strategy forcing failure against
an adversary: it constructs one possible bad execution.

\section{Worker I: guarded affine resources}
\label{sec:resource}
\status{Implemented by both Python producers; checked by a separately written arithmetic/partition replay.}

\subsection{The accepted transition language}
Fix a finite $Q$ and dimension $d$. States are $(q,x)$ with $x\in\N^d$, ordered
componentwise at equal control. A transition is
\begin{equation}
 t=(q,q',a,A,b),\qquad
 x\geq a,\qquad x'=A(x-a)+b,
 \label{eq:transition}
\end{equation}
where $a,b\in\N^d$ and $A\in\N^{d\times d}$. All assignments are simultaneous
and read the old state. The guard is explicit; it is not dropped because
natural subtraction happens to be total in Lean.

Ordinary token transitions use $A=I$. A reset is a zero row; a transfer or
merge puts several unit entries in one row; duplication permits entries
larger than one. Different finite locations encode different control states.
The prototype deliberately rejects negative entries, floating-point
coefficients and unsupported fields, including zero tests.

\begin{lemma}[Monotonicity]
\label{lem:resource-mono}
Every transition \eqref{eq:transition} is strongly upward compatible.
\end{lemma}
\begin{proof}
If $x\leq y$ and $x\geq a$, then $y\geq a$ and
$y-a=(x-a)+(y-x)$. Since $A$ is nonnegative,
$A(y-a)+b=A(x-a)+b+A(y-x)\geq A(x-a)+b$. Both executions use the same
transition and end at the same control.
\end{proof}

Dickson's finite-product property makes $\N^d$ WQO. A finite union of its
copies, with equality on the control component, is WQO as well. This supplies
the termination hypothesis, but the replay checker needs only the local
arithmetic statements in this section and the finite certificate theorem.
The relevant finite-product and antichain formalization infrastructure is
already present in Mathlib.\cite{mathlibwqo}

\subsection{Reducing an infinite preimage to a finite integer box}
For a target vector $u$ at $q'$, put
\begin{equation}
 r_j=\max(u_j-b_j,0),\qquad z=x-a.
\end{equation}
The exact predecessor condition is
\begin{equation}
 x\geq a\quad\text{and}\quad Az\geq r.
 \label{eq:preineq}
\end{equation}
For each input coordinate define
\begin{equation}
 C_i=\max_{j:A_{ji}>0}\left\lceil\frac{r_j}{A_{ji}}\right\rceil,
 \qquad \max\varnothing=0.
 \label{eq:cap}
\end{equation}

\begin{lemma}[Coordinate clipping]
\label{lem:clipping}
If $z\in\N^d$ satisfies $Az\geq r$, then
$\bar z_i=\min(z_i,C_i)$ also satisfies $A\bar z\geq r$.
More generally, this holds for any natural $C_i$ satisfying
$A_{ji}C_i\geq r_j$ for every positive entry $A_{ji}$.
\end{lemma}
\begin{proof}
Fix a row $j$. If no coordinate contributing positively to that row was
clipped, its sum is unchanged. Otherwise choose one clipped coordinate $i$
with $A_{ji}>0$. Its contribution after clipping is
$A_{ji}C_i\geq r_j$, already sufficient to meet the whole row demand. All
other terms are nonnegative. Thus every row remains satisfied, even when
several coordinates are clipped simultaneously.
\end{proof}

\begin{corollary}[Exact finite predecessor basis]
\label{cor:prebasis}
Let $F=\{z\in\N^d:0\leq z\leq C,\ Az\geq r\}$. Then
\begin{equation}
 P_t(u)=\{a+z:z\in\Min F\}
 \quad\text{satisfies}\quad
 \Pre_t(\Up{\{(q',u)\}})=\Up{\{(q,p):p\in P_t(u)\}}.
\end{equation}
\end{corollary}
\begin{proof}
Every displayed predecessor is feasible. Conversely clip any feasible
residual $z$ into the finite box. Among feasible points below the clipped
point, choose a componentwise minimal one. Adding $a$ gives a listed
predecessor below the original state. Minimal solutions cannot lie outside
the box, since clipping would produce a strictly smaller feasible point.
\end{proof}

This bound is elementary but particularly useful at the verification boundary:
the checker verifies multiplications $A_{ji}C_i\geq r_j$. It need not repeat
the producer's ceiling/max implementation. Zero rows with positive demand
make the feasible set empty, which is also detected by the replay.

\subsection{A cheap exact rule for diagonal maps}
If $A$ is diagonal, each coordinate is independent. For $A_{ii}>0$ the
least residual is $\lceil r_i/A_{ii}\rceil$. For $A_{ii}=0$, a positive
$r_i$ makes the predecessor set empty, and $r_i=0$ imposes no residual
requirement. There is at most one minimal predecessor. This includes ordinary
Petri transitions:
\begin{equation}
 \operatorname{pred}_t(u)=a+\max(u-b,0)
 \qquad(A=I).
 \label{eq:petri-pre}
\end{equation}
The implementation takes this branch before considering any box enumeration.
A huge numeric threshold therefore need not create a huge local search grid.

\subsection{Independent finite-box replay}
The eager producer enumerates $F$ and discards dominated residual vectors.
The checker does \emph{not} rerun that enumeration. It receives the proposed
minimal residuals, clipping bounds, and a binary integer-box partition.
For an inclusive box $[L,H]$, the permitted leaves are:
\begin{description}[style=nextline,leftmargin=1.5em]
\item[Covered.] A proposed feasible residual $m$ satisfies $m\leq L$.
Every point in the box lies above it.
\item[Impossible.] A row $j$ satisfies $\sum_i A_{ji}H_i<r_j$.
By nonnegativity, every point in the box violates that row.
\item[Split.] For some $i$ and $L_i\leq k<H_i$, split into the exact disjoint
integer boxes obtained by replacing $H_i$ by $k$ and $L_i$ by $k+1$.
\end{description}
The checker also tests that the proposed residuals are feasible, in the box,
sorted and pairwise incomparable. Coverage of the whole root box and clipping
then establish \cref{cor:prebasis}. A redundant or malformed list cannot be
accepted merely because its elements happen to be feasible.

The scheme is complete as a finite certificate language for these boxes.
Repeated splitting reaches singletons. A feasible singleton is above a minimal
feasible residual; an infeasible singleton violates some row. This is a
mathematical completeness fact, not a promise that an enormous tree fits the
program's resource limit.

\subsection{Worked example: all semaphore populations at once}
Write the resource state as $(n,c,\ell)$: idle processes, critical processes,
and available locks. The transitions are
\[
 \mathsf{enter}: (n,c,\ell)\mapsto(n-1,c+1,\ell-1),\quad n,\ell\geq1,
\]
\[
 \mathsf{leave}: (n,c,\ell)\mapsto(n+1,c-1,\ell+1),\quad c\geq1.
\]
The bad target is $c\geq2$, represented by $(0,2,0)$. The executed producer
returns exactly
\begin{equation}
 B=\{(0,2,0),\ (1,1,1),\ (2,0,2)\}.
 \label{eq:sembasis}
\end{equation}
The two additional states arise by taking predecessors of the bad target
through one and two enters. Leave predecessors introduce no additional
minimal unsafe states. The full certificate checks all six
transition/basis pairs, not just those used to discover the three nodes.

In particular, $(N,0,1)$ is safe for \emph{every} $N\in\N$: it dominates none
of \eqref{eq:sembasis}. The second and first vectors require $c\geq1$ or
$c\geq2$, while the third requires two locks. The final arithmetic proof can
be discharged in Lean after the exact-region theorem has been imported; it
contains no reachability reasoning.

A deliberately faulty enter returns the consumed lock immediately. With the
same leave transition, the computed basis is
\[
 \{(0,2,0),\ (1,1,0),\ (2,0,1)\}.
\]
The middle state really is unsafe even with no available lock: leave first
creates a lock, after which two faulty enters reach the error. This is a useful
example of why the final threshold set should be computed from \emph{all}
transitions rather than guessed from the visibly broken instruction. The
prototype exports a concrete execution for each unsafe demonstration query.

\subsection{Merging and amplification}
For a transition $(x,y)\mapsto(x+y,0)$ to an error-checking location with
target $(4,0)$, the source basis is
\[
 (0,4),(1,3),(2,2),(3,1),(4,0).
\]
Add a source self-loop $x\mapsto2x$ guarded by $x\geq1$, leaving $y$
unchanged. The source unsafe basis shrinks to just
\[
 (0,4),\quad(1,0).
\]
The implementation retains ten derivation nodes while the final basis across
both controls has three states. Deleting obsolete antichain members must not
delete the historical reasons still referenced by the surviving nodes.

The one-dimensional fixture $x\mapsto2x$ with the same guard and target
$x\geq10^{100}$ reaches basis $\{1\}$ after 334 admitted/expanded nodes.
The extracted execution from $x=1$ contains 333 doubling steps. The producer
does not execute $10^{100}$ steps and makes no such claim: it uses exact
large integers and logarithmically many backward thresholds.

\section{Worker II: lossy FIFO channels}
\label{sec:fifo}
\status{Implemented for a finite alphabet, a fixed finite number of channels, and single-symbol send/receive/tau transitions.}

\subsection{The order is subsequence, not prefix or substring}
A channel contains a finite word over a finite alphabet $\Sigma$. Write
$u\sqsubseteq w$ when $u$ is obtained by deleting zero or more symbols from
$w$, preserving order. Thus $ab\sqsubseteq aabb$ and
$ab\sqsubseteq bab$, whereas $ab\not\sqsubseteq ba$.
For $k$ channels, compare each of their words independently at equal control.
Higman's lemma and finite products make this order WQO. Mathlib's
formalization uses \code{List.SublistForall\textsubscript{2}} over a base
relation; specializing the base relation to equality gives the intended
embedding order.\cite{mathlibwqo}

The prototype's alphabet is a list of distinct single-character symbols.
This is an implementation encoding of finite letters, not a license to erase
arbitrary data payloads into characters. A Lean port should use a finite type
such as \code{Fin alphabetSize}, together with a proved source encoding.

\subsection{Explicit loss semantics}
Each action permits a loss before the reliable action and a loss afterward.
For one active channel, the equivalent direct relations are
\begin{align}
 w\xrightarrow{!a}v&\quad\Longleftrightarrow\quad v\sqsubseteq wa,\\
 w\xrightarrow{?a}v&\quad\Longleftrightarrow\quad av\sqsubseteq w,\\
 w\xrightarrow{\tau}v&\quad\Longleftrightarrow\quad v\sqsubseteq w.
 \label{eq:losssem}
\end{align}
Other channels may lose symbols but do not perform a send or receive during
this action. These relations preserve FIFO order. A receive consumes a head
symbol \emph{after} possible prefix losses; it cannot reorder or invent
messages. This agrees with standard lossy-channel semantics.\cite{flat}

Monotonicity is immediate: a larger input can first delete down to the
smaller one and reproduce its execution. Extra pure loss steps need not be
added to the transition list for upward-target coverability: the predecessor
of an upward region under pure loss is that same region.

\subsection{Exact principal predecessors}
For a target word $u$, define
\begin{equation}
 \rho_{!a}(u)=
 \begin{cases}
 v&u=va,\\
 u&\text{otherwise},
 \end{cases}
 \qquad
 \rho_{?a}(u)=au,
 \qquad
 \rho_\tau(u)=u.
 \label{eq:wordpre}
\end{equation}
Only the active channel changes; all other target words are preserved.

\begin{lemma}[Single-symbol predecessor law]
\label{lem:wordpre}
For each of the three action kinds,
\[
 \Pre_\theta(\Up{\{u\}})=\Up{\{\rho_\theta(u)\}}.
\]
\end{lemma}
\begin{proof}
For receive, a possible post-state covering $u$ exists exactly when
$au\sqsubseteq w$. Necessity follows by composing the embeddings through
the received state. Sufficiency deletes to $au$, consumes its first symbol,
and leaves $u$.

For send, the condition is $u\sqsubseteq wa$. If an embedding uses the
appended last position, that position must represent the last symbol of $u$,
which must therefore be $a$. Removing it leaves the prefix $v$ embedded in
$w$. If $u=va$ but an embedding does not use the appended symbol, $u$ and
hence $v$ already embed in $w$. Conversely $v\sqsubseteq w$ gives
$va\sqsubseteq wa$. When $u$ does not end in $a$, no embedding can use the
appended position and the condition is simply $u\sqsubseteq w$. The empty
word falls into this latter case and requires no input.

For tau, loss can leave a word covering $u$ exactly when $u\sqsubseteq w$.
The same argument applies componentwise to the unchanged channels.
\end{proof}

Every local predecessor set therefore has a single minimal vector of words.
The global antichain can nevertheless contain many incomparable vectors and
can grow expensive. Neither Higman's lemma nor the one-predecessor rule gives
a small uniform runtime bound.

\subsection{A send loop and two ordered receives}
Consider controls $q_0,q_1,q_2$ with $q_2$ bad and actions
\[
 q_0\xrightarrow{!b}q_0,\qquad
 q_0\xrightarrow{?a}q_1,\qquad
 q_1\xrightarrow{?b}q_2.
\]
The computed basis is
\[
 (q_0,a),\quad(q_1,b),\quad(q_2,\eps).
\]
Without the send loop the $q_0$ basis would require $ab$. With the loop,
any queue containing an $a$ as a subsequence can append a $b$, discard
unneeded letters and perform the two receives. The certificate has four
historical nodes because the earlier $ab$ requirement is superseded by $a$.
This is unbounded queue reasoning: the result applies to all queue lengths,
not just the lengths used in the finite differential tests.

The two-channel fixture similarly records a requirement of an $a$ in one
channel and a $b$ in the other, respecting the transition order through the
finite controls. Componentwise comparisons never let surplus data on one
channel satisfy a requirement on the other.

\subsection{Lossy is not reliable}
The model with one transition $q_0\xrightarrow{?a}q_1$, with $q_1$ bad,
is unsafe from queue $ba$ under lossy semantics. Delete $b$, then receive
$a$. Under reliable semantics that receive is blocked and there is no bad
execution in this two-control example.

Consequently a lossy counterexample is not automatically a counterexample
to reliable FIFO code. Reliable executions embed in lossy executions, so
\emph{safety} of the over-approximation can transfer to reliable code once
the source bridge is proved. A reported lossy failure remains an abstract
witness unless it passes a separate reliable replay. The provided
\path{lossy_not_reliable} fixture makes this distinction observable.

\section{Do not enumerate a local basis just to prove it is covered}
\label{sec:separation}
\status{A second producer and an additional independently checked receipt language are implemented and tested.}

\subsection{The avoidable intermediate object}
The eager algorithm constructs $P_t(u)$ and then checks that each of its
members lies in the current global region $\Up B$. But the necessary closure
obligation is only
\begin{equation}
 \Pre_t(\Up{\{u\}})\subseteq\Up B.
 \label{eq:directclosure}
\end{equation}
A full minimal predecessor basis is sufficient evidence, not logically
necessary evidence. This difference can be enormous.

Suppose a transition consumes $e_1$, merges $d$ residual counters into one,
and is asked to reach threshold $h$. Its minimal residuals satisfy
$z_1+\cdots+z_d=h$, so there are
\begin{equation}
 \binom{h+d-1}{d-1}
 \label{eq:mincount}
\end{equation}
of them. If the current global source basis already contains $e_1$, every
predecessor is covered \emph{solely by the guard}. Constructing all of
\eqref{eq:mincount} before checking coverage is wasted work.

\subsection{Direct global-region partition receipts}
Use the same clipping box $[0,C]$ as in \cref{lem:clipping}. In place of local
minimal residuals, permit a covered leaf to refer directly to a final global
basis state $v$ at the source control, and require
\begin{equation}
 v\leq a+L.
 \label{eq:globalleaf}
\end{equation}
Impossible and split leaves are unchanged. Coverage of the root box proves
that every feasible residual in that box yields a state in $\Up B$.
For an arbitrary feasible residual $z$, clipping gives
$\bar z\leq z$ and a covered $a+\bar z$; upward closure then covers $a+z$.
Thus the tree proves \eqref{eq:directclosure} directly.

The immutable derivation DAG still supplies the opposite inclusion. Omitting
the local Pareto basis weakens neither the final exactness statement nor the
counterexample authority. It merely removes an intermediate data structure
that the upper-inclusion theorem did not require.

\subsection{Closure separation as a search primitive}
A local procedure now has two possible successful answers:
\begin{equation}
 \mathsf{Covered}(\text{partition proof})
 \quad\text{or}\quad
 \mathsf{Uncovered}(p,\text{concrete predecessor step}).
 \label{eq:separation}
\end{equation}
The latter point need not be a \emph{locally} minimal predecessor. It only
needs to be feasible and outside the currently represented upward region.

The implemented recursive box search follows this order:
\begin{lstlisting}
probe(L, H):
    if some global basis v satisfies v <= consume + L:
        return a global-basis coverage leaf
    if some matrix row is below its demand at H:
        return an impossible leaf
    if L itself meets every row demand:
        return the uncovered state consume + L
    split the widest non-singleton coordinate at its midpoint
    recursively inspect both children
\end{lstlisting}

At a feasible singleton not covered by the region, the third branch must
return a witness. At an infeasible singleton, the second branch must succeed.
Therefore the procedure terminates on a finite box and returns either a
sound coverage tree or an actual uncovered predecessor, absent cutoffs.

On receiving a witness, the producer admits its derivation node, updates the
antichain, and probes again. It continues until the transition is covered,
or until the target node itself is dominated and will be superseded by an
enqueued smaller target. It processes every surviving target and every
incoming transition. At the end it reconstructs all final coverage trees
against the final basis indices: an old tree cannot be reused with stale
references after the antichain changes.

\begin{theorem}[Correctness of separation-based saturation]
\label{thm:separation}
With exact local separation, strong upward compatibility, a WQO, and no
operational cutoffs, the separation-based worklist terminates and produces the
same minimal unsafe basis as the eager worklist.
\end{theorem}
\begin{proof}
Every admitted state is an uncovered predecessor of an existing derivation
node, hence has a valid lower-inclusion reason and strictly enlarges the
represented upward set. Infinite admission is excluded by the same bad-sequence
argument as in \cref{thm:termination}. Once admissions stop, each remaining
probe terminates with coverage. Surviving nodes receive all required coverage
checks, giving \eqref{eq:closurecover}; target coverage was preserved.
Apply \cref{thm:exact}. Since the final basis is an antichain in a partial
order, it is the unique set of minimal unsafe states, independently of which
uncovered predecessors were chosen.
\end{proof}

\subsection{Executed ablation, with the exact claim}
The two producers were run on the same 190 models and produced identical
final bases. The separation checker independently replayed those 190
certificates and three additional dense-preimage fixtures. These are repeated
models, not 380 distinct benchmark systems.

\begin{table}[htbp]
\centering\small
\begin{tabular}{@{}rrrrr@{}}
\toprule
$d$ & $h$ & Local minima by \eqref{eq:mincount} & Local grid size & Certificate bytes\\
\midrule
4&20&1,771&194,481&776\\
6&60&8,259,888&51,520,374,361&898\\
8&100&26,075,972,546&10,828,567,056,280,801&1,063\\
\bottomrule
\end{tabular}
\caption{Constructed globally covered preimages. Minima/grid counts are exact formula evaluations, not objects enumerated during the experiment. Bytes count the full canonical compact-JSON region certificate.}
\label{tab:ablation}
\end{table}

Each fixture has two source-to-target transitions. The first consumes $e_1$
and produces at least the bad threshold, establishing the source basis
$e_1$. The second is the guarded dense merger described above. The eager
producer returns \code{unknown} at its 100,000-point local-grid limit. The
separation producer returns a complete certificate with two derivation nodes,
two final basis states, and a one-leaf direct proof for the dense preimage.
Its two counted partition visits are one during search and one during final
certificate construction. No giant basis was generated and no giant trace
was expanded.

This is an ablation of \emph{our two Python implementations}, designed to
expose an unnecessary intermediate representation. It is not evidence that
existing optimized Petri-net tools fail on the fixtures, and it is not a
comparison against Lean or \code{grind}.

\subsection{Further concrete optimizations, not yet implemented}
\paragraph{Factor the arithmetic constraint graph.}
Form a bipartite graph with one node for each residual variable, one for each
positive-demand row, and an edge for each positive coefficient. Disconnected
components involve disjoint variable sets. Handle isolated rows immediately:
a positive-demand row with no incident variable is impossible. Solve each
remaining component independently. The local feasible set is a Cartesian
product; its minimal vectors are exactly products of componentwise minima.
This can reduce arithmetic work greatly, although explicitly listing the
product may still be exponentially large. A direct global-cover tree can
avoid that product when global cones suffice.

\paragraph{Use an integer solver only to find uncovered points.}
An optional Z3 search can assert $z\geq0$, $Az\geq r$, and
\[
 \bigwedge_{v\in B\text{ at source}}
       \bigvee_i (a_i+z_i<v_i).
\]
This is quantifier-free linear integer arithmetic because $A$ is fixed.
A satisfying model proposes an uncovered predecessor, which is independently
checked. Coordinatewise minimization within that model's lower box can use
successive bounded integer queries. The accepted local point need not be
optimal. An \code{unsat} answer is \emph{not} the closure proof: request or
construct the direct partition receipt and replay it. Z3's documented
\code{SolverFor("QF\_LIA")}, integer terms and models provide a concrete
backend for this optional search path.\cite{z3}

\paragraph{Index dominance without changing its meaning.}
Store resource states in a separate dominance index per control. A trie over
coordinates with subtree lower/upper bounds can reject whole subtrees when
a prefix cannot dominate the query; a final exact componentwise comparison
still decides admission. For words, length and letter-count vectors are cheap
\emph{necessary} embedding filters, followed by an exact subsequence test.
Equal lengths or equal Parikh vectors are not sufficient for embedding.
Caching these filters is an optimization, not a quotient of the state space.

\section{What Forge should export after checking}
\label{sec:output}

\subsection{An exact weakest safety precondition}
For resource states at control $q$, the exact result is the finite formula
\begin{align}
 \Unsafe(q,x)
   &\Longleftrightarrow
     \bigvee_{b\in B_q}\ \bigwedge_{i=1}^{d}(b_i\leq x_i),\\
 \Safe(q,x)
   &\Longleftrightarrow
     \bigwedge_{b\in B_q}\ \bigvee_{i=1}^{d}(x_i<b_i).
 \label{eq:cnf}
\end{align}
Empty disjunction and conjunction conventions give the correct answers for
an empty bad set. A basis vector of all zeros makes its whole control unsafe.
The meaning is the weakest initial condition for avoiding the specified bad
set under every finite execution of the exact model.

For a symbolic family $x=I(\theta)$, substitute that expression in
\eqref{eq:cnf}. For affine natural initializers the residual formula is
Presburger arithmetic, so Forge can route it to its ordinary arithmetic
layer. Polynomial or other initializers can use existing specialized workers
without changing the region certificate. The new engine resolves the
unbounded transition closure; the existing engine resolves the remaining
initial-condition implication.

\subsection{Executable witnesses, not just basis membership}
An unsafe initial state is assigned to any basis node below it. Follow the
node's child links. For resource transitions, execute the original matrix
update on the actual current state. Do \emph{not} replace that state by the
smaller recorded child: the larger values can affect subsequent execution.
Monotonicity proves that the original action word remains enabled.

For lossy transitions, emit the exact pre-loss words, the result of the
reliable action, and the post-loss words. The independent trace checker
verifies each deletion direction, that receives consume the current head,
that sends append at the end, and that the final state covers an original
bad target. It does not trust a table of claimed embeddings or an output
label saying ``counterexample.''

The implementation offers two distinct artifacts: the region certificate and
the concrete trace for a requested initial state. A valid trace can refute
safety for that exact model and initial state even without a region
certificate. A valid region certificate classifies every initial state even
when no concrete trace has yet been requested.

\section{A concrete Lean integration plan}
\label{sec:lean}
\status{Specified, not implemented or compiled in this work. The new Python checkers are not formally verified.}

\subsection{Keep the trust target small and semantic}
Forge already specifies its acceptance boundary and has some elaborated
core-only specimens; this work does not undo that progress or claim that
nothing in Forge has compiled. What is missing \emph{here} is a Lean
implementation and soundness proof of these particular checkers, followed by
an original-goal bridge.\cite{forgestatus}

The first deliverable in Lean should be a \emph{finite-certificate theorem},
not a formalization of a complicated terminating search. The search may remain
in Python, or later move to a faster language. The checker needs no WQO theorem
to justify an accepted finite certificate. WQO formalization can be a separate
completeness milestone. This separation makes it possible to get a useful
proof-producing worker without first formalizing all search scheduling,
dominance indexing or solver internals.

Use Forge's existing runtime and closure infrastructure, adding modules with
roughly the following responsibilities. The names below are proposed project
paths, not modules that currently exist.
\begin{longtable}{@{}p{0.36\textwidth}p{0.59\textwidth}@{}}
\toprule
Proposed module & New obligation it owns\\
\midrule
\endhead
\path{Forge/Ordered/Region.lean} & Upward regions, finite derivation DAG interpretation, complete transition/basis coverage, exact-region theorem.\\
\path{Forge/Ordered/Resource.lean} & Guarded nonnegative-affine semantics, strong monotonicity, residual conversion, clipping and diagonal predecessor lemmas.\\
\path{Forge/Ordered/BoxCover.lean} & Integer partition interpretation and both local-minimum and direct-global-cover soundness.\\
\path{Forge/Ordered/Lossy.lean} & Word/channel semantics, exact send/receive/tau predecessor laws and concrete loss replay.\\
\path{Forge/Ordered/Decode.lean} & Versioned data decoding into bounded, dimension-correct internal structures. No semantic authority in raw JSON.\\
\path{Forge/Ordered/Reify.lean} & Source-to-model equivalence or explicitly one-way simulation; original target correspondence.\\
\path{Forge/Ordered/Replay.lean} & Evaluate the checker, build its correctness application, return the exact original proposition.\\
\path{Forge/Ordered/Complete.lean} & Optional formalized WQO termination and equality of eager/separation mathematical algorithms.\\
\bottomrule
\end{longtable}

\subsection{A theorem dependency graph}
The new statements should be developed in the following order.
\begin{enumerate}[leftmargin=2em]
\item Prove that a finite sequence of compatible transitions lifts upward.
This is the only generic monotonicity lemma needed by DAG replay.
\item Prove DAG-node soundness by induction on the node index; prove backward
closure soundness by induction on finite executions. Combine them into
\cref{thm:exact}.
\item For resources, prove \cref{lem:resource-mono}, the guarded residual
identity, and \cref{lem:clipping}. Then prove each partition constructor
sound, and conclude each predecessor receipt sound.
\item For queues, prove the operational equivalences in \eqref{eq:losssem},
then \cref{lem:wordpre}, then the concrete pre/post-loss replay theorem.
\item Prove that the executable checker accepting data implies all semantic
premises of the region theorem. Finally connect the decoded model to the
source relation and bad predicate.
\end{enumerate}

All numeric dimensions should be explicit in Lean. A resource vector can have
type \code{Fin d $\to$ Nat}; a matrix can have type
\code{Fin d $\to$ Fin d $\to$ Nat}; control values can have type
\code{Fin qCount}. A word can be a \code{List (Fin alphabetSize)}.
The executable storage may be arrays or lists for efficiency, but a length
proof must connect it to this semantic interpretation. No out-of-range index
should become an implicit default zero.

The final theorem should express equality or pointwise equivalence of the
unsafe set, not just a Boolean classification for an unbound serialized
state. The intended mathematical signature is
\[
 \mathsf{check}(M,C)=\mathsf{true}
 \ \Longrightarrow\
 \forall s,\quad
 (\exists z,\ \Reach_M(s,z)\land\Bad_M(z))
 \ \Longleftrightarrow\ s\in\Up{B_C}.
\]
The exact Lean declaration names and syntax remain implementation work; this
article does not ship an uncompiled \code{sorry}-filled file pretending to
satisfy the contract.

\subsection{Source-model bridges: three different permission levels}
Let the Lean source relation be $\Rightarrow$ on $X$, its bad predicate be
$D$, and an encoder be $E:X\to S$.

\paragraph{Exact encoding.}
For a representation equivalence, require correspondence of initial states,
bad predicates, and transitions, with enough back-translation to reconstruct
source executions. For a noninjective quotient, a bare existential equality
of quotient edges is insufficient: abstract steps must lift from \emph{every}
represented concrete state where the result is claimed. Otherwise independently
realizable abstract edges can form a spurious path.

One adequate contract has source forward simulation, target equivalence
$D(x)\leftrightarrow\Bad(E(x))$, and a step-lifting property
\[
 E(x)\to s'\Longrightarrow
 \exists x',\ x\Rightarrow x'\land E(x')=s'.
\]
With these properties, exact model classification transfers to the source.
More flexible relational simulations are possible, but their lifting
quantifiers must be written just as carefully.

\paragraph{Safety over-approximation.}
For a safety-only bridge it is enough that source steps simulate to model
steps (or finite model paths), and $D(x)\Rightarrow\Bad(E(x))$.
Model safety then implies source safety. Model unsafety is only a candidate
counterexample, because it may exploit extra behaviors. The reliable-to-lossy
queue translation is the central example.

\paragraph{Witness-only under-approximation.}
A concretely replayed source execution establishes a source failure even when
no exact or over-approximate source model is available. This gives no authority
to declare unobserved states safe. The three entry points should remain
separate in the API, rather than sharing a Boolean \code{verified} flag.
Forge's existing evidence-boundary design supplies this discipline; these
are its new, model-specific proof obligations.\cite{forgeideas}

\subsection{An implementable resource reifier}
Start with an explicit-model entrance rather than silently analyzing arbitrary
Lean code. The user supplies a finite transition list and a theorem connecting
it to the source relation. This immediately supports serious proofs and
isolates the checker work from general program extraction.

The next entrance can recognize inductively defined finite transition
relations whose rules normalize to lower-bound guards and simultaneous
nonnegative-affine updates. It should enumerate every constructor, assign a
stable source reference, extract $a,A,b$, and prove the transition equation
under the original guard. Polynomial normalization can help establish the
coefficient identity, but ordinary natural subtraction must be rewritten only
under the relevant lower-bound hypotheses. If a rule has an upper-bound guard,
a zero test, a negative coefficient, an unbounded control variable or a
value-dependent matrix, refuse the exact entrance rather than erase it.

A source expression such as $2x-2$ is not globally interchangeable over
naturals with an unconstrained integer expression. Under $x\geq1$ it equals
$2(x-1)$, and that conditional equality is what the model bridge must prove.
Likewise a sequence of assignments must first be converted to a simultaneous
old-state expression; using the matrix of the written right-hand sides before
substitution changes the program.

For anonymous finite-state populations, a counting quotient is especially
valuable. Concrete states are multisets of agents; each transition consumes
and produces specified multiplicities. A sufficient-count theorem chooses the
required distinct agents and proves that every abstract step lifts. This
justifies exact classification for every population size. Identity-sensitive
protocols, nonexchangeable payloads and hidden local state need a stronger
quotient proof or a safety-only abstraction, not an automatic exact label.

\subsection{An implementable queue reifier}
The initial queue entrance should accept an explicit finite-control relation
over lists with single-symbol sends/receives and explicit lossy semantics.
Provide the reliable-to-lossy safety bridge as a separate theorem. Enumerate
and check every channel, control, symbol and action; an omitted action would
make the computed safe set unsound for the original model.

For data-carrying messages, a finite alphabet map needs a proved behavioral
contract. A many-to-one map may be an over-approximation, but it need not
support counterexample transfer. The original source's equality tests,
branching on payloads and typeclass dictionaries must remain accounted for.
In particular, this worker does not solve the Church-indexing synthesis gap
reported in the current Leant/Djex diagnostics. Applying subsequence pruning
to arbitrary synthesized functions would be unsound without the very
representation and behavior theorems that those diagnostics carefully keep
separate.\cite{leant,djex,indexdiag}

\subsection{Concrete libraries and their roles}
\begin{table}[htbp]
\centering\small
\begin{tabularx}{\textwidth}{@{}p{0.25\textwidth}XX@{}}
\toprule
Library/tool & Proposed use & Trust/implementation status\\
\midrule
Lean 4 and Mathlib & Semantic statements, finite induction, integer arithmetic, lists, WQO infrastructure & Required for kernel acceptance; port not executed here\\
Python standard library & Exact integer prototypes, JSON receipts, deletion-mask oracles & Executed; no third-party runtime dependency\\
Z3 integer API & Optional uncovered-point proposal and local minimization & Not integrated or used by the delivered producers; solver answers are not receipts\\
TAPAAL / verifypn & Optional untimed token-net corpus and external producer adapter & Not integrated; admit only a proved supported subset, not arbitrary timed/game/inhibitor semantics\\
LoLA & Optional external token-net comparison and concrete trace source & Not integrated; independent source/model replay still required\\
\bottomrule
\end{tabularx}
\caption{Specific tools with explicit boundaries, rather than a list of solvers presumed to produce Lean proofs.}
\end{table}

Mathlib's \path{Mathlib/Order/WellFoundedSet.lean} contains the finite-product
and antichain infrastructure, including
\path{Set.PartiallyWellOrderedOn.pi} and
\path{IsAntichain.finite_of_partiallyWellOrderedOn}. Its documented Higman
entry point is
\path{Set.PartiallyWellOrderedOn.partiallyWellOrderedOn_sublistForall₂}.
\path{Mathlib/Order/WellQuasiOrder.lean} supplies the underlying WQO vocabulary.
These names were checked against the online documentation; the exact imports
and elaboration at Forge's pinned dependency still need a compiler run.
The pipeline must not infer that a documentation lookup is an API acceptance
test.\cite{mathlibwqo,mathlibbasic}

Forge's reviewed baseline specifies Lean 4.34.0 and Mathlib commit
\path{1cf325a0cf67aca2b04d76b5380ff6a9e410aefa}. Use that pair for the first
integration branch instead of mixing the repository with a moving Mathlib
head.\cite{forge} TAPAAL and LoLA are concrete external verification projects,
not certified Lean backends out of the box. Their adapters are optional future
work.\cite{tapaal,lola}

\section{Executed experiments and what they establish}
\label{sec:experiments}

\subsection{Environment and reproducibility}
The recorded run used Python 3.13.5 on x86\_64 Linux and seed 941572603.
All delivered search and replay code runs under \code{python -S}, disabling
site packages. The test corpus, raw accepted certificates, concrete
illustrative traces, summaries and run output are included. Timings are
single-run wall-clock observations, stored for diagnosis, not controlled
performance estimates.

The corpus generator, explicit oracle and arithmetic checker are separate
functions/modules. The checker imports neither producer. In a fresh replay
process, an import hook makes attempts to import the search module or
NumPy/SciPy/SymPy raise an exception. This verifies a real execution-path
separation. It does not prove that independently written routines lack a
shared mathematical misunderstanding.

\begin{table}[htbp]
\centering\small
\begin{tabularx}{\textwidth}{@{}Xrr@{}}
\toprule
Experiment and unit & Cases/queries & Discrepancies\\
\midrule
Local resource predecessor problems & 1,296 & 0\\
Resource-point membership comparisons on those problems & 63,504 & 0\\
FIFO action/target/input word comparisons & 2,325 & 0\\
Complete finite-resource execution queries, 100 systems & 7,000 & 0\\
Complete acyclic FIFO execution queries, 40 systems & 2,040 & 0\\
Complete length-nonincreasing FIFO queries, 40 systems & 2,480 & 0\\
Eager region certificates independently replayed & 190 & 0 rejected\\
Same-model separation/eager basis comparisons & 190 & 0\\
Separation region certificates independently replayed & 193 & 0 rejected\\
Eager/general parser and semantic corruption attempts & 32 & 0 accepted invalid\\
Additional direct-cover corruption attempts & 5 & 0 accepted invalid\\
Operational-cutoff controls & 4 & 0 promoted to a verdict\\
\bottomrule
\end{tabularx}
\caption{Distinct units remain distinct. Several rows intentionally reuse the same models; there is no combined success-rate denominator.}
\label{tab:tests}
\end{table}

\subsection{Local predecessor tests}
For the resource lane, the generator covers every $2\times2$ matrix with
entries in $\{0,1,2\}$ and every target in $\{0,1,2,3\}^2$: $81\cdot16=1296$
problems. Consumption/production vectors are chosen deterministically by the
recorded seed. On each problem, it compares the represented predecessor set
with direct forward arithmetic for all 49 points of $\{0,\ldots,6\}^2$.
This includes points beyond the clipping bounds and points below the guards.
The finite input family is not exhaustive over arbitrary dimensions,
coefficients or guards.

For FIFO, there are five concrete action variants: send/receive either binary
letter, plus tau. Every binary input word of length at most four is compared
against every target word of length at most three. The oracle enumerates
actual deletion masks before and after the reliable action. It does not call
\eqref{eq:wordpre}, the producer's greedy subsequence predicate, or the
checker's dynamic-programming embedding predicate. The $5\cdot31\cdot15=2325$
comparisons all agree.

\subsection{Why the finite execution oracles really terminate}
The resource systems have two controls and three counters. Every matrix
column has sum at most one, and total production is at most total consumption.
For an enabled transition,
\[
 \|A(x-a)+b\|_1\leq\|x-a\|_1+\|b\|_1\leq\|x\|_1.
\]
Each queried initial state has total resource at most four. The full reachable
space lies in a finite simplex, so explicit BFS exhausts the actual reachable
states rather than truncating an infinite search. The 7,000 queries divide
into 910 unsafe and 6,090 safe results, with exact agreement.

For the acyclic channel family, finite control paths bound the number of
sends; initial total length is at most two and there are two channels. For
the cyclic family there are no sends, so total length never increases. These
are independent reasons why explicit execution terminates. The acyclic
queries divide into 124 unsafe and 1,916 safe; the decreasing-length queries
into 1,319 unsafe and 1,161 safe. The oracle has a defensive state limit, but
hitting it would raise an inconclusive test failure; it never becomes a safety
answer. No query in the recorded run reached that limit.

The unbounded send-loop and amplification examples are covered by their
finite certificate proofs, not by claiming that short explicit tests enumerate
all their executions. This distinction is important even when every displayed
test is green.

\subsection{Corruption and refusal controls}
The tests remove required closure rows and target coverage; duplicate rows
and basis entries; create a cyclic child reference; corrupt a lower-inclusion
state; tamper with clipping bounds, minima and partition nodes; incorrectly
select the diagonal rule; invent messages during pre/post losses; and try
invalid numeric values and duplicate JSON keys. All these intentionally
invalid objects are rejected.

A source transition is also changed and the certificate subject deliberately
rebound to that changed model. Rejection in this case depends on the semantic
checks, not merely a mismatching subject string. Unsupported zero-test fields
are rejected rather than ignored. The additional separation tests corrupt
its direct global-basis references, partition and row completeness.

The four producer cutoff controls independently exhaust the node, expansion,
local-grid and box-certificate budgets. All return \code{unknown}. These tests
exercise selected failure paths; they do not establish resistance against
hostile or arbitrarily large input. The decoder has explicit byte, dimension,
word-length, item-count and integer-bit limits, but it is research code rather
than a security-audited sandbox.

\subsection{A preserved failed run}
The first experiment run reached the named-example assertions and failed
because the reporting dictionary used the key \code{basis} both for the
list of basis states and for a numeric statistic. Merging the statistics
replaced the list with its length. The produced semaphore certificate already
contained the expected three vectors. The repair nested the statistics under
\code{search\_stats}; the experiment then ran again. The original traceback
and run output are retained in \path{results/failed-runs/}. This was a
reporting/test-harness defect, not evidence of a false region certificate;
its scope should not be exaggerated in either direction.

\subsection{Unmeasured claims}
No theorem in this new package was checked by Lean. No \code{forge} tactic was
installed. No Mathlib-dependent source was compiled. No benchmark against
\code{grind}, Aesop, \code{omega}, \code{nlinarith}, TAPAAL, LoLA or an external
SMT solver was executed. The old repository's test counts are not added to
these. The article's mathematical proofs, the Python results and a future
Lean-kernel acceptance are three distinct levels of evidence.

\section{Implementation priorities and acceptance gates}
\label{sec:roadmap}

\subsection{First merge: a bounded, explicitly invoked worker}
The first useful integration should accept an explicit resource model and
produce the exact-region theorem for diagonal transitions. This covers
ordinary token systems, resets on independent coordinates and guarded
amplification, while requiring no general box proof. Acceptance requires a
compiled soundness theorem, successful replay of the original certificate
against the pinned Lean/Mathlib pair, an original-source semantic bridge,
and rejection of the deliberately faulty version. A theorem about a
hand-rewritten model is not enough unless its equivalence to the source is
also proved.

Next add the resource clipping lemma and direct-global-cover checker. This
is more useful than first formalizing exhaustive local Pareto enumeration:
it admits both producers' useful results while avoiding the representation
explosion demonstrated in \cref{tab:ablation}. Keep the eager local-minimum
receipt as a second checker path and an independent test fixture.

Then add the lossy-channel semantics and predecessor laws. The first acceptance
corpus should include the send-loop example, the two-channel example, a
correct concrete loss trace, and the reliable/lossy distinction. A test suite
that validates only positive lossy examples has not checked the source
abstraction boundary.

\subsection{Required regression tests for the Lean port}
Beyond porting the stored arithmetic and word cases, add exact source tests
for simultaneous versus sequential updates, guarded natural subtraction,
missing transition constructors, control reindexing, mismatched channel
identities and many-to-one payload encodings. Test an upward target and a
nearby exact target together, ensuring the latter is refused or carries an
explicit approximation theorem.

For certificate construction, check the original goal in a clean verification
context according to Forge's existing policy. Check the axiom dependencies of
the resulting declarations and the source bridge. A theorem that assumes
\code{checkSound} as an axiom is not a certified checker. An elaborated
structure containing fields named \code{sound} is not a proof that any
serialized input satisfies those fields.

For performance evaluation, separate source reification, producer time,
certificate size, Lean replay time, and final residual-goal solving. Compare
with baseline tactics on the same original source statements and with the same
provided lemmas. Do not count a difficult model lemma as free preparation for
one side while charging its proof to the other. A worker may be valuable even
when its examples also have short manual invariants; the principal new metric
is coverage of whole parametric initial-state families without supplied
invariants.

\subsection{Boundaries that should remain visible}
The resource lane does not support zero tests or arbitrary negative updates.
The queue lane does not decide reliable FIFO reachability, exact-word targets,
infinite alphabets or fairness. The current implementation does not combine
resource counters and queues in one model, although a product-domain worker
with appropriate independent action semantics is a plausible next extension.
It does not solve mixed games or probabilistic reachability by changing the
interpretation of ``unsafe.'' Those subjects already have different workers
and proof contracts in Forge.

WQO guarantees eventual termination of the idealized supported algorithms,
not practical runtime. Large antichains and long bad sequences remain genuine
costs. The direct-cover language avoids some intermediate explosions; it does
not remove all output-size or state-space complexity. Timeout, memory refusal,
unsupported reification and lack of a source bridge remain legitimate,
explicit nonacceptance outcomes.

\section{Conclusion}
Forge should gain the ability to return an exact finite unsafe-region basis
for two new infinite-state languages: guarded nonnegative-affine resources
and lossy FIFO channels. The same result yields parameterized safety
conditions, minimal resource/message requirements for failure, and a finite
family of concrete action plans for every unsafe initial state.

The proposal is implementable in stages because finite certificate soundness
does not depend on verifying the search's WQO termination proof. Its most
important representation decision is to let a checker accept direct coverage
by the final global antichain rather than insist on enumerating every local
minimal predecessor. The executed separation experiment illustrates the
consequence with a 1,063-byte certificate where the unnecessary local
predecessor basis has over 26 billion elements.

What exists now is a detailed semantic design, proved mathematical reductions,
two executed Python producers, an independent finite-data checker, replayable
witnesses and a reproducible test archive. What remains is the Lean checker and
the original-source bridge. Those are precise engineering and formalization
tasks, not gaps to be hidden behind a claim that the prototype already is a
Lean tactic.

\appendix
\section{Certificate format and executable entry points}
\label{app:format}

\subsection{Region object}
The versioned JSON object has the following top-level fields:
\begin{lstlisting}
{
  "format": "forge-ordered-v1",
  "subject": "canonical encoding of the complete original model",
  "nodes": [ ... immutable earlier-child derivations ... ],
  "basis": [ ... node indices in canonical state order ... ],
  "target_cover": [ ... one basis index per original bad target ... ],
  "closure": [
    { "basis": 0,
      "transition": 2,
      "predecessor": { ... receipt ... },
      "cover": [ ... basis indices ... ] }
  ]
}
\end{lstlisting}
The subject is the complete canonical model encoding, not a short hash. The
checker receives the expected model separately and compares this encoding.
The check is not a source-goal fingerprint: that extra relationship belongs
to the future Lean reifier and source bridge.

Each derivation is either an exact original-target leaf or a state with a
transition index and strictly smaller child index. The checker rejects unknown
fields and out-of-range indices. Sorted states use control then lexicographic
data order solely for deterministic encoding; semantic dominance is the
separate componentwise or subsequence partial order.

\subsection{Local receipt tags}
The \code{diagonal} tag requires an actually diagonal matrix and uses exact
integer quotient/remainder rules. The \code{box} tag contains natural clipping
caps, feasible minimal residuals and a local-minimum partition. The
\code{lossy-word} tag triggers the proved send/receive/tau predecessor rule.
For these three tags, the closure row carries one coverage index per computed
predecessor.

The \code{resource-cover} tag contains clipping caps and a partition whose
covered leaves refer directly to final global basis indices. Its row carries
an empty local \code{cover} list. All incoming transition/basis pairs must
still occur exactly once. This direct tag does not bypass the lower-inclusion
DAG or original-target coverage checks.

A future schema revision should preserve explicit version negotiation. It
should not let an unknown receipt tag fall through to a default semantics, or
let a malformed basis index mean the zero vector.

\subsection{Commands}
From the archive root:
\begin{lstlisting}
# Reproduce the full first suite into a new directory.
python -S prototype/run_tests.py --output reproduced-results

# Compare the second producer on the same generated models.
python -S prototype/run_separation.py \
  --source reproduced-results --output reproduced-separation

# Independently check a stored example; no search is invoked.
python -S prototype/checker.py \
  results/examples/semaphore.problem.json \
  results/examples/semaphore.certificate.json

# Search and independently check a new explicit model.
python -S prototype/cli.py \
  results/examples/semaphore.problem.json --out new-region.json

# Replay the largest dense-preimage receipt.
python -S prototype/checker.py \
  results/separation/dense-8.problem.json \
  results/separation/dense-8.certificate.json
\end{lstlisting}
The CLI defaults to the separation producer. Its successful output explicitly
says \code{accepted\_by\_python\_checker} and
\code{lean\_kernel\_checked: false}. It refuses to overwrite an existing
certificate path. Optional initial-state input exports an independently
checked concrete trace when the state is unsafe.

To rebuild this article, run \code{pdflatex} on
\path{article/forge_ordered.tex} from its directory three times. No external
bibliography processor, shell escape, network connection, or third-party
Python library is required by the artifact.

\section{Evidence inventory}
\begin{longtable}{@{}p{0.43\textwidth}p{0.52\textwidth}@{}}
\toprule
Artifact & Meaning\\
\midrule
\endhead
\path{prototype/search.py} & Both untrusted producers; exact arithmetic, subsequence search, witness extraction.\\
\path{prototype/checker.py} & Independently written region and concrete-trace replay; no producer import.\\
\path{prototype/run_tests.py} & Deterministic local tests, complete finite oracles, corruption and cutoff controls.\\
\path{prototype/run_separation.py} & Same-model producer ablation and dense-preimage fixtures.\\
\path{prototype/cli.py} & Explicit-model command line; checker acceptance required before a success output.\\
\path{results/summary.json} & Environment, grouped counts, named basis states, queries and raw single-run timing observations.\\
\path{results/corpus.json} & 190 model/certificate pairs from the eager run.\\
\path{results/trace-corpus.json} & 18 illustrative concrete traces from named examples; not all oracle witness checks are stored here.\\
\path{results/separation/corpus.json} & 193 model/certificate pairs: the same 190 models plus three dense-preimage fixtures.\\
\path{results/separation/summary.json} & Exact-basis agreement, constructed-case size calculations and direct-cover mutation results.\\
\path{results/failed-runs/} & Preserved first reporting-harness failure and its original output.\\
\path{docs/novelty-audit.md} & The review boundary, reused ideas, delta, and non-exhaustiveness caveat.\\
\path{docs/LEAN_PORT_PLAN.md} & Missing formalization work and explicit acceptance conditions.\\
\bottomrule
\end{longtable}

\clearpage
\begingroup
\small
\begin{thebibliography}{99}
\bibitem{forge}
V. Reshetnikov, \emph{Forge}, README at commit
\path{58ea206bd7ad501b930add568151217bcc7f82a2}.
\url{https://github.com/VladimirReshetnikov/Forge/blob/58ea206bd7ad501b930add568151217bcc7f82a2/README.md}.

\bibitem{forgestatus}
V. Reshetnikov, \emph{Forge implementation status}, same snapshot.
\url{https://github.com/VladimirReshetnikov/Forge/blob/58ea206bd7ad501b930add568151217bcc7f82a2/docs/STATUS.md}.

\bibitem{forgeclosure}
V. Reshetnikov, \emph{Finite certificates for unbounded behaviour}, Forge
merged article, Section~6 source, same snapshot.
\url{https://github.com/VladimirReshetnikov/Forge/blob/58ea206bd7ad501b930add568151217bcc7f82a2/article/sections/06-closure.tex}.

\bibitem{forgeideas}
V. Reshetnikov, \emph{Ideas from Leant and Djex}, Forge, same snapshot.
\url{https://github.com/VladimirReshetnikov/Forge/blob/58ea206bd7ad501b930add568151217bcc7f82a2/docs/IDEAS-FROM-LEANT-DJEX.md}.

\bibitem{leant}
V. Reshetnikov, \emph{Leant}, README at commit
\path{6bf05ad78c467989e68290f2d08bbed40802d485}.
\url{https://github.com/VladimirReshetnikov/Leant/blob/6bf05ad78c467989e68290f2d08bbed40802d485/README.md}.

\bibitem{djex}
V. Reshetnikov, \emph{Djex}, commit record
\path{e8778f4ebd63e1f9b9b410fa4de8d14a8a04c9e5}, ``Locate the remaining indexing gap in bounded composition search.''
\url{https://github.com/VladimirReshetnikov/Djex/commit/e8778f4ebd63e1f9b9b410fa4de8d14a8a04c9e5}.

\bibitem{indexdiag}
V. Reshetnikov, \emph{Indexing composition after native kind adoption},
14 September 2026, Leant diagnostic report at the reviewed snapshot.
\url{https://github.com/VladimirReshetnikov/Leant/blob/6bf05ad78c467989e68290f2d08bbed40802d485/docs/reports/2026-09-14-indexing-composition-diagnostics.md}.

\bibitem{acjt}
P.~A. Abdulla, K. \v{C}er\=ans, B. Jonsson and Y.-K. Tsay,
``Algorithmic analysis of programs with well quasi-ordered domains,''
\emph{Information and Computation} 160(1--2), 109--127, 2000.
\url{https://doi.org/10.1006/inco.1999.2843}.

\bibitem{fs}
A. Finkel and P. Schnoebelen, ``Well-structured transition systems everywhere!,''
\emph{Theoretical Computer Science} 256(1--2), 63--92, 2001.
\url{https://doi.org/10.1016/S0304-3975(00)00102-X}.

\bibitem{flat}
P. Schnoebelen, \emph{On flat lossy channel machines}, arXiv:2007.05269v1,
2020. In particular, Sections~2 and~3.1 for loss semantics and predecessors.
\url{https://arxiv.org/html/2007.05269v1}.

\bibitem{mathlibwqo}
The Mathlib community, \emph{Mathlib.Order.WellFoundedSet}, online
API documentation, inspected 15 September 2026.
\url{https://leanprover-community.github.io/mathlib4_docs/Mathlib/Order/WellFoundedSet.html}.

\bibitem{mathlibbasic}
The Mathlib community, \emph{Mathlib.Order.WellQuasiOrder}, online
API documentation, inspected 15 September 2026.
\url{https://leanprover-community.github.io/mathlib4_docs/Mathlib/Order/WellQuasiOrder.html}.

\bibitem{z3}
The Z3 project, \emph{Python API reference}, integer expressions, solvers,
models and \code{SolverFor}.
\url{https://z3prover.github.io/api/html/z3.z3.html}.

\bibitem{tapaal}
The TAPAAL project, official project site and verification tools.
\url{https://www.tapaal.net/}.
\url{https://github.com/TAPAAL/verifypn}.

\bibitem{lola}
The LoLA project, official tool page, University of Rostock.
\url{https://theo.informatik.uni-rostock.de/en/theo-forschung/tools/lola/}.
\end{thebibliography}
\endgroup
\end{document}
